% !TEX spellcheck = de
\DocumentMetadata{lang=de}
\documentclass{ltx-talk}

\makeatletter
\providecommand*{\input@path}{}
\g@addto@macro\input@path{{./include/}{./beamertheme-sam/}}
\makeatother

\definecolor{theme}{HTML}{870030}

\usepackage[
  themecolor=theme,
]{talkthemesam}

\AddToHook{section/begin}{
  \UseTemplate{sectionpage}{talk}{}
}

\colorlet{lightgray}{samlgray!20}
\colorlet{red}{themecolor}
\tcbset{lefthand ratio=0.65}

\usepackage[ngerman]{babel}
\usepackage{csquotes}
\usepackage{lucide-icons}
\usepackage{tabularray}
\UseTblrLibrary{booktabs}
\UseTblrLibrary{tikz}
\UseTblrLibrary{counter}
\UseTblrLibrary{amsmath}

\usepackage[style=authoryear]{biblatex}
\addbibresource{\jobname.bib}
\DeclareFieldFormat{doi}{%
  \newline
  \mkbibacro{DOI}\addcolon\space
  \ifhyperref
    {\href{https://doi.org/#1}{\nolinkurl{#1}}}
    {\nolinkurl{#1}}}
\DeclareFieldFormat{url}{\newline\mkbibacro{URL}\addcolon\space\url{#1}}    

\title{Tabellen mit Tabularray}
\titlegraphic{Schloss_Biebrich,_Nordseite,_Wiesbaden-Biebrich,_150904,_ako}
\sectiongraphic{section}
\titlecredit{
  Titelbild: \href{https://commons.wikimedia.org/wiki/File:Schloss_Biebrich,_Nordseite,_Wiesbaden-Biebrich,_150904,_ako.jpg}{Ansgar Koreng (CC BY 3.0 via wikimedia commons)}%
}
\author{samcarter}
\institute{Dante Herbsttagung Wiesbaden}
\date{25.\,-- 26.\ September 2026}

\begin{document}

\maketitle

% Was ist tabularray
\begin{frame}
  \frametitle{Was ist Tabularray}
  \begin{columns}
    \begin{column}{.65\textwidth}
      \begin{itemize}
        \item modernes Paket für Tabellen und Arrays
        \item veröffentlicht in 2021 von Jianrui Lyu
        \item adoptiert von den TeXackers
        \item erlaubt Trennung von Format und Inhalt
        \item sehr mächtig (heute Grundlagen und Highlights)
      \end{itemize}   
    \end{column}
    \begin{column}{.35\textwidth}
      \hfill\color{themecolor}\lucideicon[height=4cm]{message-circle-question-mark}\hspace*{1cm}
    \end{column}
  \end{columns}  
\end{frame}

% einfaches Beispiel
\begin{frame*}
  \frametitle{Einfaches Beispiel}
  \begin{tcblisting}{}
    \usepackage{tabularray}
    
    \begin{tblr}{}
      a & b & c\\
      1 & 2 & 3\\
      4 & 5 & 6\\      
    \end{tblr}
  \end{tcblisting}  
\end{frame*}

% Viele Wegen führen nach Rom
\begin{frame*}
  \frametitle{Viele Wege führen nach Rom...}
  \tcbset{equal height group=diabox, valign=center}%
  \vskip-0.4cm
  
  \begin{columns}[vertical-alignment=top]
    \begin{column}{.3\textwidth}
      \begin{tcblisting}{
        sidebyside=false,
        listing above text,
        title=Direkt,
        before lower=\medskip\centering
      }
        \begin{tblr}{}
          \hline
          a & b & c\\
          1 & 2 & 3\\
          \hline
        \end{tblr}
      \end{tcblisting}  
    \end{column}
    \begin{column}{.3\textwidth}
      \begin{tcblisting}{
        sidebyside=false,
        listing above text,
        title=Tabellen-Präamble,
        before lower=\medskip\centering
      }
        \begin{tblr}{
          hline{1,Z}={},
        }
          a & b & c\\
          1 & 2 & 3\\ 
        \end{tblr}
      \end{tcblisting}  
    \end{column}  
    \begin{column}{.3\textwidth}
      \begin{tcblisting}{
        sidebyside=false,
        listing above text,
        title=Global,
        before lower=\medskip\centering
      }
        \SetTblrInner{
          hline{1,Z}={},
        }
        \begin{tblr}{}
          a & b & c\\
          1 & 2 & 3\\ 
        \end{tblr}
      \end{tcblisting}  
    \end{column}        
  \end{columns}
\end{frame*}

\section{Werkzeugkiste}

%farben
%
%zellen können noch mehr, mode text imath, cmd
%
% columns, rows
%
%alignment, mergen, hspan
%
%typen von tablllen, longtblr, talltblr

% childindex

% Linien
\begin{frame*}
  \frametitle{Horizontale und vertikale Linien}
  \begin{tcblisting}{}
    \begin{tblr}{
      hline{1}={solid},
      hline{2}={2-Z}{dashed,fg=red,wd=2pt},
      vlines
    }
      a & b & c\\
      1 & 2 & 3\\
      4 & 5 & 6\\ 
    \end{tblr}
  \end{tcblisting}  
\end{frame*}

% mehrzeilige Texte
\begin{frame*}
  \frametitle{Mehrzeilige Texte}
  \begin{tcblisting}{}
    \begin{tblr}{vlines,hlines}
      a & {c \\ d}\\
    \end{tblr}
  \end{tcblisting}  
\end{frame*}

% zählen
\begin{frame*}
  \frametitle{Das kleine Einmaleins}
  \begin{tcblisting}{}
    \begin{tblr}{}
      x & \therownum & x \\
      \therownum & x & x \\
      \thecolcount/\therowcount & x & x \\
    \end{tblr}
  \end{tcblisting}
  \begin{tcblisting}{}
    \begin{tblr}{
      row{1}={cmd=\thecolnum},
      column{1}={cmd=\therownum},
    }
      &   &  \\
      & x & x\\
    \end{tblr}
  \end{tcblisting}  
\end{frame*}

\section{Bibliotheken}

% Mathe
\begin{frame*}
  \frametitle{Mathematik-Umgebungen mit amsmath}
  \begin{tcblisting}{}
    \UseTblrLibrary{amsmath}
    
    \begin{equation*}
      \begin{+pmatrix}[
        cell{2}{2}={bg=lightgray}
      ]
        a & b & c\\
        1 & 2 & 3\\
        4 & 5 & 6\\ 
      \end{+pmatrix}
    \end{equation*}
  \end{tcblisting}
  \bigskip
  
  Weitere: \saminline|+array|, \saminline|+matrix|, \saminline|+bmatrix|, \saminline|+Bmatrix|, \saminline|+vmatrix|, \saminline|+Vmatrix| und \saminline|+cases|
\end{frame*}

% TikZ
\begin{frame*}
  \frametitle{Ti\emph{k}Z}
  \begin{tcblisting}{}
    \UseTblrLibrary{tikz}
    
    \begin{tblrtikzabove}
      \draw[red,ultra thick,->] 
        (table.north west) -- (2-2.south east);
    \end{tblrtikzabove}%
    \begin{tblr}{}
      a & b & c\\
      1 & 2 & 3\\
      4 & 5 & 6\\ 
    \end{tblr}
  \end{tcblisting}
\end{frame*}

% booktabs
\begin{frame*}
  \frametitle{Professionelle Tabellen mit booktabs}
  \begin{tcblisting}{}
    \UseTblrLibrary{booktabs}
   
    \begin{tblr}{}
      \toprule
      a & \SetCell[c=2]{halign=c} b & \\
      \cmidrule[r]{1}\cmidrule[l]{2-3}
      c & d & e\\
      \midrule
      1 & 2 & 3\\
      4 & 5 & 6\\
      \bottomrule 
    \end{tblr}
  \end{tcblisting}
\end{frame*}

% weitere Bibliotheken
\begin{frame*}
  \frametitle{Weitere Bibliotheken:}
  \begin{columns}
    \begin{column}{.65\textwidth}
      \begin{itemize}
        \item siunitx
        \item varwidth
        \item hook
        \item counter
        \item babel (via \saminline|\usepackage{tblr-extras}|)
        \item caption (via \saminline|\usepackage{tblr-extras}|)
      \end{itemize}
    \end{column}
    \begin{column}{.35\textwidth}
      \begin{itemize}
        \item diagbox
        \item html
        \item nameref
        \item zref
        \item functional      
      \end{itemize}
    \end{column}  
  \end{columns}
\end{frame*}

% Nachteile
\begin{frame}
  \frametitle{Nachteile}
  \begin{columns}
    \begin{column}{.65\textwidth}
      \begin{itemize}
        \item langsamer als klassische Tabellen
        \item keine getaggted Tabellen
        \item umständliche Expansion/optionale Argumente
      \end{itemize}   
    \end{column}
    \begin{column}{.35\textwidth}
      \hfill\color{themecolor}\lucideicon[height=4cm]{message-circle-warning}\hspace*{1cm}
    \end{column}
  \end{columns}
\end{frame}

% Noch neugierig
\begin{frame}
  \frametitle{Noch neugierig?}
  \begin{itemize}
    \item \fullcite{tabularray}
    \item \fullcite{github}
    \item \fullcite{Maggi:2023:DBP}
  \end{itemize}
\end{frame}

\end{document}